%*-------------------------------------------------------------------------*
%*                 Copyright 2013 Core Physics, Inc.                       *
%*  Under terms of the contract to support CASL, there is a non-exclusive  *
%*  license for use of this work by or on behalf of the U.S. Government.   *
%*-------------------------------------------------------------------------*
%
% @version CVS $Id: materials.tex,v 1.11 2013/11/19 17:28:04 scott Exp $
%
%%%%%%%%%%%%%%%%%%%%%%%%%%%%%%%%%%%%%%%%%
\chapter{Materials Description}  \label{chap:materials}

This chapter contains a description of the material input.  There are two types of materials
in the input file -- structural materials (input with a {\it mat} card) and fuel materials
(input with a {\it fuel} card).

Structural materials can be defined in either the [CORE] block, or in the
geometry object blocks [ASSEMBLY], [INSERT], [CONTROL], and [DETECTOR].
If the materials are defined in the [CORE] block, they have global scope.
If the materials are defined in the geometry object blocks, they only have
scope in the block they are defined.
The reason for this is to maintain the modularity of the geometry objects.

Fuel materials can only be defined in [ASSEMBLY] blocks.

Materials are used in many different input cards.  They are used to define
cells, nozzles, core plates, baffles, grids, reflectors, etc.  Every material that
is used must be defined with either a {\it mat} card or a {\it fuel} card.

%%%%%%%%%%%%%%%%%%%%%%%%%%%%%%%%
\section{Structural Materials}

Structural materials are not fuel and do not deplete.  Structural
materials are defined with the following input card:

{\bf mat} {\it user-mat} {\it density} ({\it library-name$_i$}, {\it frac$_i$}, i=1, I)

where:
\begin{itemize}
\item {\it user-mat} is a user-defined material name.  It is case sensitive.  {\it user-mat} is used to
define material names in other input cards such as {\it cell}, {\it grid}, {\it nozzle}, etc. (No default)
\item {\it density} is the material density in g/cc (No default)
\item {\it library-name} is a corresponding material library name for the user material.  The library name
must be defined in the material composition file. (Default = {\it user-mat}).  Multiple library materials
can be mixed to form a single user material.
\item {\it frac}  is the fraction of the library material in the user material. (Default=1.0 if there is only one library material in the user material).
\end{itemize}

There are two special user materials,  ``mod'' and ``vacuum''.
The user can use these materials in cell definitions, but the code will be automatically
determine the composition of these materials  based on T/H feedback and soluble boron concentrations.
The user is not allowed to define a user materials named ``mod'' or ``vacuum'' on a {\it mat} card.

Some example material cards are shown below.
\begin{verbatim}
  mat zirc 6.56                ! use default library name and fraction
  mat zirc 6.56 zirc4 1.0
  mat B10  12.0 boron 1.0
  mat XYZ  6.0  zirc4 0.8 ss 0.2
  mat ABCD 8.0  zirc4 0.8 ss 0.15 b4c 0.05
\end{verbatim}

All of the material fractions must sum to either +1.0 or -1.0.
If positive fractions are used, the fractions refer to atomic fractions.
If negative fractions are used, the fractions refer to weight fractions.


%%%%%%%%%%%%%%%%%%%%%%
\subsection{Search Order}

Structural materials can be defined in either the [CORE] block or one of the
geometry object blocks.  When a material is referred to in a block, it will look
for the material definition in the following order:
\begin{enumerate}  % enumerate = numbered list
 \item The code will first look for the material name in the local block ([ASSEMBLY], [INSERT], [CONTROL], or [DETECTOR])
 \item If the material is not found in the local block, it will look in the [CORE] block
\end{enumerate}

If materials are defined in the [CORE] block, they have global scope over the entire input.
If materials are defined in other blocks, they only have scope over the local block.  This means
that two geometry object blocks can use different material definitions with the same name.
One example of this is that two assemblies can be defined with the material ``zirc'', but
the ``zirc'' can have different compositions in the blocks.

%%%%%%%%%%%%%%%%%%%%%%%%%%
\section{Composition File}

All library materials must be included in the HDF composition file.
The name of the composition file is specified on the [INSILICO] {\it mat\_library} card,
or the [MPACT] {\it (TBD)} card, and is usually named something like ``casl\_comp.sh5''.

A partial list of materials that are present on the current composition file is:
\begin{itemize}   % itemize = bullet list
  \item aic
  \item b4c
  \item boron
  \item inc
  \item he
  \item mod
  \item ss
  \item zirc
\end{itemize}


%%%%%%%%%%%%%%%%%%%%%%%%
\section{Fuel Materials}

Fuel materials are composed mostly of UO$_2$ with different U-235 enrichments.
Fuel materials are different from structural materials in that they deplete and have
additional properties as described below.

Fuel can only be defined in [ASSEMBLY] blocks, and fuel materials can only be referenced by {\it cell} cards
in the [ASSEMBLY] block they are defined in.

Fuel materials are defined with the following input card:

{\bf fuel} {\it user-mat} {\it density} {\it thden} / {\it U-235\_enrichment} \{{\it HM\_material$_i$}={\it HM\_enrichment$_i$}, i=1, N\}  \\
   \{ / {\it gad\_material}={\it gad\_fraction} \}

Where:

\begin{itemize}   % itemize = bullet list
  \item {\it user-mat} is a user-defined fuel name. It is case sensitive. (No default)
  \item {\it density} is the fuel material density in g/cc (No default)
  \item {\it thden} is the percent of theoretical density in the pellet (\%) (No default)
  \item {\it U-235\_enrichment} is the U-235 enrichment in the fuel in weight \%  (No default).
  \begin{itemize}   % itemize = bullet list
    \item If U-234 and U-236 are not specified, the enrichments will automatically be added to the
       fuel by a pre-determined function (see below)
    \item If the sum of the heavy metal (HM) enrichments does not equal 100\%, the remainder of the
        HM composition will be assigned to U-238.
  \end{itemize}
  \item {\it HM\_material$_i$} is the material name for a HM (Pu-239, Pu-241, etc.) (optional)
  \item {\it HM\_enrichment$_i$} is the enrichment of HM isotope $i$ in weight \% (optional)
  \item {\it gad\_material} is the material name for the gad oxide (or other material) (optional).
        The gad material may be a mixture defined on a separate {\it mat} card.
  \item {\it gad\_fraction} is the weight percent of the gad material relative to the total mass (optional)
\end{itemize}

Oxygen should not be included on the {\it fuel} card.  The correct amount of oxygen will automatically
be added to the HM to create an oxide (either UO$_2$  or (HM)O$_2$).

The {\it density} is the ``stack density'' or ``smeared density'' and should include the volume
of the pellet dishing and chamfers.
It is calculated as the total mass of the fuel pellets divided by the total volume of the fuel
\begin{equation}
  \mbox{stack density} = \frac{ (\mbox{fuel mass})} { \pi (\mbox{pellet radius})^2 \, (\mbox{fuel height}) }
\end{equation}

The {\it thden} refers to the actual density of the pellet.  This quantity is
only used in fuel performance codes to evaluate material properties.

If U-234 or U-236 enrichments are not included, they are automatically added to the fuel with the following formulas:
\begin{equation}
   W_{234} = 0.007731 \cdot W_{235}^{1.0837}
\end{equation}
\begin{equation}
  W_{236} = 0.0046 \cdot W_{235}
\end{equation}
Where $W_{23x}$ is the enrichment of each of the uranium isotopes in percent.


Examples of valid fuel cards are:
\begin{verbatim}
  fuel U21 10.4 95.2 / 2.1
  fuel UO2-35 10.297 95.0 / 3.5
  fuel U23 10.111 / 2.3                      ! bare minimum input for fuel
  mat gad5 7.407 gd2o3 1.0                   ! define gad material
  fuel U49 10.111 94.5 / 1.8  / gad5=0.05    ! 5% gad
  fuel PX1 10.111 94.5 / 0.711  Pu-239=23.4  Pu-241=1.9    ! Sample MOX
  fuel UPX21G 10.111 94.5 / 2.1 U-234=0.001 Pu-239=3.4  Pu-241=1.9 / gad3=0.03
\end{verbatim}



